% ============================================================================
%  Section VI — Results and Discussion  (IEEE TIM, IEEEtran)
%  Passive voice / third person throughout.
%  Requires: booktabs, graphicx, pifont.
%  Figures (results/): fig_optimism_gap.png, fig_confusion_group.png,
%    fig_recipe_stacking.png, fig_tradeoffs_main.png.
%  Numbers from results/analysis_C_multiseed.*, analysis_D.*, analysis_harden.*,
%    analysis_E.*.
\providecommand{\cmark}{\ding{51}}
\providecommand{\xmark}{\ding{55}}
% ============================================================================

\section{Results and Discussion}
\label{sec:results}

\subsection{The Optimism Gap of Random-Split Evaluation}
\label{subsec:gap}
Table~\ref{tab:optimism} reports the classification accuracy of the three
baseline models under the four evaluation protocols, and Fig.~\ref{fig:gap}
visualises the same result. Under the RANDOM protocol every model attains
essentially perfect accuracy, reproducing---and slightly exceeding---the
$98\%$ reported by the source data article's random-split baseline. When the same
features and models are evaluated under the leakage-free GROUP protocol, accuracy
collapses to roughly one-half, an \emph{optimism gap} of the order of $40$
accuracy points. The CROSS-SPEED and CROSS-LOAD protocols, which additionally
impose an operating-condition shift, yield comparably reduced accuracy. The
near-perfect RANDOM result is therefore an artefact of evaluation rather than
evidence of diagnostic capability: because overlapping windows share half of their
samples and every window of a recording shares one operating point and sensor
coupling, a randomly partitioned model can identify the originating recording
instead of the fault. This finding is consistent with the leakage effects
reported for industrial-grade vibration data~\cite{wheat2024impact,
hendriks2022towards, kapoor2023leakage} and demonstrates that the concern applies
equally, and with large magnitude, to low-rate smartphone-based acquisition.

Two further observations are of practical relevance. First, the leakage-free
accuracy, although far below the random-split figure, remains well above the
$16.7\%$ chance level, confirming that genuine fault information is recoverable
from $100$~Hz smartphone data. Second, the lower-capacity models generalise at
least as well as the Random Forest under the GROUP protocol, indicating that the
forest overfits the small number of training recordings.

The seed-averaged confusion matrix of the GROUP protocol
(Fig.~\ref{fig:confusion}) explains where the residual difficulty lies. The
healthy state and the broken-rotor-bar fault are the most reliably identified,
whereas the two lubrication-severity states (B1 and B2) are frequently confused
with each other and with the remaining bearing conditions. This pattern is
consistent with the sensor characterisation of Section~\ref{sec:charact}: the
distinguishing signatures of incipient bearing degradation reside in the
high-frequency band that the $100$~Hz sampling cannot represent, so coarse states
are separable while fine bearing-severity staging is not.

\subsection{Accuracy--Resource Trade-offs}
\label{subsec:tradeoff}
The trade-off between diagnostic accuracy and the sensing and communication cost
of an edge node was characterised by varying one operating parameter at a time
under the GROUP protocol. Counter-intuitively, reducing the sampling rate did not
degrade accuracy: as shown in Table~\ref{tab:sampling}, every rate at or below
$25$~Hz matched or exceeded the accuracy obtained at $100$~Hz, with a maximum near
$25$~Hz. The anti-alias decimation removes high-frequency content that the model
otherwise overfits, while the discriminative low-frequency information identified
in Section~\ref{sec:charact} is retained. Reducing the channel set was similarly
beneficial: the three raw axes alone outperformed the full six-channel set,
consistent with the channel-informativeness ordering established earlier, whereas
lengthening the analysis window improved accuracy at the cost of detection
latency.

Because these frugal choices are individually beneficial, their combination was
examined (Table~\ref{tab:recipe} and Fig.~\ref{fig:recipe}). Applying the
resource-optimal configuration---$25$~Hz sampling, the three raw channels, and a
$4$~s window---raised the Random Forest accuracy from $53.7\%$ at the naive
$100$~Hz/six-channel/$2$~s baseline to $74.2\%$, while reducing the raw data
volume by approximately eight-fold. The resource-optimal node is therefore not
merely adequate but \emph{superior} to the naive configuration under honest
evaluation. Robustness to further deployment stresses was also confirmed:
reducing the amplitude resolution to $4$~bits lowered accuracy by only a few
points, injected packet loss of up to $20\%$ left accuracy essentially unchanged,
and the six-feature subset of the source baseline performed within noise of the
full feature set. The principal sensitivity was to sensor mounting: a simulated
$30^{\circ}$ change in orientation degraded accuracy substantially, indicating
that consistent placement, or orientation-invariant features, is required in
practice.

\subsection{Communication and Edge-Node Feasibility}
\label{subsec:edge}
The implications of the trade-off analysis for an IoT deployment were quantified
through a communication link budget and an edge-model footprint.
Table~\ref{tab:linkbudget} compares the bitrate of three telemetry strategies
against the sustainable throughput of common low-power radios. Streaming the raw
six-channel signal requires $9.6$~kbps and transmitting the extracted feature
vector requires $2.69$~kbps; both are readily carried by Bluetooth Low Energy,
Zigbee, or NB-IoT, but neither fits within the duty-cycle-limited throughput of a
LoRaWAN uplink, which admits only the $40$~bps on-node decision stream. For
LoRa-class links, therefore, classification must be executed on the node.

Such on-node execution is feasible with the present pipeline.
Table~\ref{tab:footprint} lists the leakage-free accuracy and the serialised size
of four candidate models. A logistic-regression classifier of $4.7$~kB matched or
exceeded the accuracy of a $1.09$~MB Random Forest, and the time-domain feature
extraction requires only $\mathcal{O}(WC)$ operations per window with no spectral
transform. Compact, edge-deployable models are thus also the better generalisers,
a convenient alignment for low-cost nodes. A cycle-accurate estimation of
inference latency and energy on a representative microcontroller is identified as
future work.

\subsection{Limitations}
\label{subsec:limits}
Several limitations bound the scope of these findings. The dataset comprises a
single $1.1$~kW motor with laboratory-induced, discrete-severity faults, and the
two lubrication states were graded qualitatively; the reported accuracies are
therefore a deployment-honest baseline rather than a claim of field performance.
The small number of recordings ($36$) produces wide confidence intervals under
the recording-holdout protocols, and although five-seed repetition stabilises the
central estimates, larger multi-machine datasets are needed to narrow them. The
edge-node analysis quantifies communication load and model footprint but reports
energy only indirectly; direct on-device measurement was outside the present
scope. Finally, generalisation across machines was not assessed, and constitutes
the most important direction for future validation.

% ---------------------------------------------------------------------------
\begin{table}[t]
\centering
\caption{Accuracy (\%) by evaluation protocol and model
(mean $\pm$ std.\ over cross-validation folds; $100$~Hz, six channels, $2$~s).}
\label{tab:optimism}
\renewcommand{\arraystretch}{1.2}
\begin{tabular}{@{}lccc@{}}
\toprule
\textbf{Protocol} & \textbf{RF} & \textbf{KNN} & \textbf{SVM} \\
\midrule
RANDOM (leaky)            & $100.0$ & $100.0$ & $100.0$ \\
GROUP (recording-holdout) & $47.6\pm7.9$ & $63.3\pm11.8$ & $53.2\pm12.2$ \\
CROSS-SPEED               & $49.4\pm11.2$ & $48.2\pm9.1$ & $49.7\pm12.9$ \\
CROSS-LOAD                & $50.1\pm4.8$ & $40.3\pm0.9$ & $40.1\pm2.0$ \\
\bottomrule
\end{tabular}
\end{table}

% ---------------------------------------------------------------------------
\begin{table}[t]
\centering
\caption{Sampling-rate sweep under the GROUP protocol (Random Forest,
five-seed mean $\pm$ std). Lower rates match or exceed $100$~Hz.}
\label{tab:sampling}
\renewcommand{\arraystretch}{1.2}
\begin{tabular}{@{}cccc@{}}
\toprule
\textbf{Rate (Hz)} & \textbf{Raw bitrate (kbps)} & \textbf{Accuracy (\%)} \\
\midrule
100  & 9.6 & 53.7 $\pm$ 13.4 \\
50   & 4.8 & 42.0 $\pm$ 12.8 \\
25   & 2.4 & \textbf{64.9 $\pm$ 15.7} \\
20   & 1.9 & 61.6 $\pm$ 16.8 \\
12.5 & 1.2 & 59.1 $\pm$ 14.9 \\
10   & 1.0 & 59.7 $\pm$ 13.7 \\
\bottomrule
\end{tabular}
\end{table}

% ---------------------------------------------------------------------------
\begin{table}[t]
\centering
\caption{Incremental frugal recipe under the GROUP protocol
(five-seed mean $\pm$ std). The resource-optimal node beats the naive baseline.}
\label{tab:recipe}
\renewcommand{\arraystretch}{1.2}
\begin{tabular}{@{}lccc@{}}
\toprule
\textbf{Configuration} & \textbf{RF} & \textbf{KNN} & \textbf{LogReg} \\
\midrule
Baseline ($100$~Hz, $6$ ch, $2$~s) & 53.7 $\pm$ 13.4 & 61.1 $\pm$ 11.7 & 57.0 $\pm$ 11.7 \\
\;$+\,25$~Hz                        & 64.9 $\pm$ 15.7 & 46.7 $\pm$ 12.9 & 59.0 $\pm$ 13.9 \\
\;$+\,$raw $3$ ch                   & 65.1 $\pm$ 15.2 & 52.8 $\pm$ 15.5 & 55.1 $\pm$ 14.8 \\
\;$+\,4$~s window (optimal)         & \textbf{74.2 $\pm$ 16.6} & 55.0 $\pm$ 19.2 & 62.6 $\pm$ 16.4 \\
\bottomrule
\end{tabular}
\end{table}

% ---------------------------------------------------------------------------
\begin{table}[t]
\centering
\caption{Communication link budget: telemetry bitrate versus sustainable uplink
capacity. \cmark~= sustainable, \xmark~= infeasible.}
\label{tab:linkbudget}
\renewcommand{\arraystretch}{1.2}
\setlength{\tabcolsep}{4pt}
\begin{tabular}{@{}lccccc@{}}
\toprule
\textbf{Telemetry stream} & \textbf{Bitrate} & \textbf{BLE} & \textbf{Zigbee} &
\textbf{NB-IoT} & \textbf{LoRaWAN} \\
\midrule
Raw ($6$ ch, int16)     & $9.6$~kbps  & \cmark & \cmark & \cmark & \xmark \\
Features ($84$/s)       & $2.69$~kbps & \cmark & \cmark & \cmark & \xmark \\
Features (min., $6$/s)  & $192$~bps   & \cmark & \cmark & \cmark & \xmark \\
On-node decision        & $40$~bps    & \cmark & \cmark & \cmark & \cmark \\
\bottomrule
\end{tabular}
\end{table}

% ---------------------------------------------------------------------------
\begin{table}[t]
\centering
\caption{Edge-model footprint versus GROUP-protocol accuracy. Compact models are
also the better generalisers.}
\label{tab:footprint}
\renewcommand{\arraystretch}{1.2}
\begin{tabular}{@{}lcc@{}}
\toprule
\textbf{Model} & \textbf{GROUP accuracy (\%)} & \textbf{Serialised size} \\
\midrule
Random Forest ($200$ trees)     & 46.4 & $1.09$~MB \\
Random Forest ($50$, depth $8$) & 51.2 & $214$~kB \\
Decision tree (depth $8$)       & 43.7 & $4.3$~kB \\
Logistic regression             & 52.2 & $4.7$~kB \\
\bottomrule
\end{tabular}
\end{table}

% ---------------------------------------------------------------------------
\begin{figure}[t]
\centering
\includegraphics[width=\columnwidth]{fig_optimism_gap.png}
\caption{Accuracy by evaluation protocol and model. Near-perfect accuracy under
the RANDOM split collapses once evaluation is leakage-free (GROUP) or imposes an
operating-condition shift (CROSS-SPEED, CROSS-LOAD).}
\label{fig:gap}
\end{figure}

% ---------------------------------------------------------------------------
\begin{figure}[t]
\centering
\includegraphics[width=\columnwidth]{fig_confusion_group.png}
\caption{Seed-averaged confusion matrix of the Random Forest under the GROUP
protocol. Healthy (H) and broken-rotor-bar (R) states are the most separable;
the lubrication-severity states (B1, B2) are the most confused.}
\label{fig:confusion}
\end{figure}

% ---------------------------------------------------------------------------
\begin{figure}[t]
\centering
\includegraphics[width=\columnwidth]{fig_recipe_stacking.png}
\caption{Incremental frugal recipe under the GROUP protocol. Combining the
resource-frugal choices raises Random Forest accuracy from $53.7\%$ to $74.2\%$
while reducing the raw data volume roughly eight-fold.}
\label{fig:recipe}
\end{figure}
